% Section 9, ethics, data and AI assistance (restructure plan, part 2). The AI-assistance statement follows the owner's
% default, "an AI coding agent", and says what the records name (the agent tool, the configurations by model tier) without
% naming a model (the final verification's finding on the model statement). The ledgers are named
% from paper/'s directory list, without a count. The misuse paragraph's clause on optimisation pressure, which points to
% the threat of that name, was added in a later round (decision-log entry D82). The privacy paragraph's sentence on the
% download records was corrected in a later round (decision-log entry D83), after a review by model agents: the committed
% index and screening manifests give every natural task's and screened comment's repository id, which the download
% records map to its URL, so the earlier sentence (from D63), which denied any such link, did not hold;
% tests/test_third_party_records.py now reads this file too. Since decision-log entry D84, which records the author
% line: the AI-assistance paragraph names a single author, and its list of what agents produced names the protocol
% and its amendments, the decision log, and this text in its drafts and revisions; the privacy paragraph says that the
% URLs stay, and that the privacy tools' defaults have no person's review.
\section{Ethics, data and AI assistance}
\label{sec:ethics}
\paragraph{Data.} The controlled sessions ran only in sandboxed copies of public, MIT-licensed repositories, and the study's agent wrote their user turns; the transcripts keep upstream authors' public names and handles as attribution, with e-mail addresses redacted. The accepted natural families' repositories are under MIT or Apache-2.0 licences, which agree with AIDev's field, and redistributed upstream code keeps its licence (notices in \texttt{THIRD\_PARTY\_LICENSES/}). No private, proprietary or customer material was used. SWE-chat was not used: its access terms were not accepted (D4).

\paragraph{Natural data and privacy.} The natural sources hold personal data such as logins and comment text. Their files stay outside the repository, and the natural tasks in a git-ignored directory with a committed index of ids, hashes and licences. The repository holds identifiers, commit SHAs, the commented files' paths, hashes of texts, counts, licences, categorical labels and checked model-written paraphrases of the qualifying review comments of both S9 samples; it holds no verbatim comment or session text, and no login or handle of the sampled pull requests' reviewers or of their pull-request, issue or commit authors (D47). The S9 and S9x download records list the repositories fetched for screening by URL, with the owning account masked as \texttt{\{owner\}} where that account acts in AIDev (whether an unmasked owner is a person's account was not checked), and the pull-request refs fetched there. They hold no task id or comment, but the committed index and screening records give every natural task's and screened comment's AIDev repository id, which these records map to that URL; these ids, URLs and refs are public, so the release links each natural task and screened comment to its repository and, on GitHub, to its pull request and its participants, which the mask does not prevent. The URLs stay: they are public and needed for provenance and for AIDev's attribution (D84). This paper prints no thread text, repository name, login or agent product of the natural data, only the study's task ids. The privacy tools' matching thresholds are builder defaults that no person has reviewed (Appendix~\ref{app:deviations}; D84).

\paragraph{Misuse.} Verifiers derived from sessions can serve as training rewards, but no policy here was optimised against a verifier, so our false-acceptance rates say nothing about such a reward under optimisation (\S\ref{sec:threats}). The process-separated verifiers resist the tampering probes we tested; B1 imports candidate code into the reporting process and is unsuitable as a reward for adversarial policies (\S\ref{sec:repro}).

\paragraph{AI assistance.} The software, the protocol and its amendments, the decision log, the scenarios and their reference operationalisation, the controlled sessions, the recovered contracts, every generated verifier (the label-support verifiers among them), the constructed candidates, the fresh attempts, the label proposals, the internal reviews and this text, in its drafts and revisions, were produced with an AI coding agent under the author's instructions. By the study's procedure its models are of one model family of one provider; the protocol and the decision log name the agent tool and define the configurations by model tier, and the paper names no model (\S\ref{sec:design-tasks}). The natural pre-review states and references are trees of AIDev pull requests, and B1 runs the repositories' own tests (in the controlled arm with the tests the session agent added). Models alone wrote and reviewed the scenarios, the reference contracts, the labels and the internal reviews: no human has written, reviewed or adjudicated any of them. Every number in the text and the tables is a generated value with an entry in an evidence ledger: \texttt{paper/evidence\_ledger.md} and its siblings \texttt{\_exploratory}, \texttt{\_fresh}, \texttt{\_natural}, \texttt{\_natural\_run}, \texttt{\_figures}, \texttt{\_pilot} and \texttt{\_design}. The author is responsible for the content.
